\documentclass[11pt,a4paper]{article}
\usepackage[T1]{fontenc}
\usepackage[utf8]{inputenc}
\usepackage{lmodern,microtype}
\usepackage[margin=25mm,headheight=15pt]{geometry}
\usepackage{amsmath,amssymb,amsthm,mathtools}
\usepackage{booktabs,longtable,array,enumitem}
\usepackage{xcolor,listings,fancyhdr}
\usepackage[hyphens]{url}
\usepackage[colorlinks=true,linkcolor=black,citecolor=black,urlcolor=blue!45!black]{hyperref}
\usepackage{cleveref}
\hypersetup{pdftitle={AsymptoticAnalysis after Six Review Waves},pdfauthor={Technical review prepared for Vladimir Reshetnikov},pdfsubject={Incremental source audit and rational Lerch certificate prototype}}
\definecolor{codeback}{gray}{0.965}
\lstset{basicstyle=\ttfamily\small,columns=fullflexible,keepspaces=true,
 breaklines=true,breakatwhitespace=false,showstringspaces=false,
 backgroundcolor=\color{codeback},frame=single,rulecolor=\color{black!15},
 xleftmargin=0.3em,xrightmargin=0.3em,aboveskip=0.9em,belowskip=0.9em}
\setlength{\emergencystretch}{3em}
\setlength{\parskip}{0.35em}
\setlength{\parindent}{0pt}
\widowpenalty=10000
\clubpenalty=10000
\setlist{nosep,leftmargin=1.5em}
\pagestyle{fancy}\fancyhf{}
\fancyhead[L]{\small AsymptoticAnalysis: incremental technical audit}
\fancyhead[R]{\small \texttt{efa1aeec}}
\fancyfoot[C]{\thepage}
\newtheorem{theorem}{Theorem}[section]
\newtheorem{proposition}[theorem]{Proposition}
\newtheorem{lemma}[theorem]{Lemma}
\theoremstyle{remark}\newtheorem{remark}[theorem]{Remark}
\newcommand{\code}[1]{\texttt{\detokenize{#1}}}
\newcommand{\R}{\mathbb R}
\newcommand{\N}{\mathbb N}
\newcommand{\PhiL}{\Phi}
\newcommand{\dd}{\,\mathrm d}
\newcommand{\repo}{\url{https://github.com/VladimirReshetnikov/Asymptotic}}
\title{\textbf{AsymptoticAnalysis after Six Review Waves}\\[0.5em]
\large Operand-preserving validation, an all-depth inverse counterexample,\\
and exact-rational Lerch certificates}
\author{Technical review prepared for Vladimir Reshetnikov}
\date{10 September 2026\\[0.6em]\small Source snapshot:\\
\texttt{efa1aeec4845a9c35e140963a0333d0c9ec33b05}}
\begin{document}
\maketitle
\thispagestyle{empty}
\begin{abstract}
This incremental audit identifies a correctness defect in the exact-core inverse
constructor: it checks target independence after adding the core and perturbation,
although it subsequently uses the operands separately. Cancellation can therefore
admit target-dependent model data into a theorem whose constants must be fixed.
For a particularly simple input, the total forward function is exactly $1/x$,
but the source-derived marker approximations alternate between $0$ and $2/y$ while
the claimed remainder improves by one power at every depth. A second family has
a genuinely vanishing perturbation ratio and still violates the reported scale.
The report gives an all-depth proof, a narrowly scoped repair, and executable
independent checks.

A separate development contribution constructs a hierarchy of rational lower and
upper bounds for $\Phi(q,1,a)$ using the two Meixner measures naturally associated
with the defining sum. The implementation evaluates two three-term recurrences:
it needs neither a moment matrix nor quadrature nodes, a special-function call,
or floating-point arithmetic. Positivity, error identities, nesting and explicit
error caps are proved. Python code is independently tested; an equivalent Wolfram
Language prototype is supplied.

The comparison baseline consists of the six review indexes, the maintained
register, and focused original passages relevant to the proposed deltas. Existing
findings are not presented as new. Evidence is deliberately separated: source
inspection, mathematical proof, isolated Wolfram 15 evaluations, and independent
Python tests do not constitute an upstream package or Mathics acceptance run.
\end{abstract}
\vfill
\textbf{Deliverable boundary.} One new source-established correctness finding;
one concrete algorithmic development proposal. No new benchmark result, complete
compatibility claim, integrated patch, or full-suite acceptance is asserted.
\newpage
\begingroup
\small
\setlength{\parskip}{0pt}
\tableofcontents
\endgroup
\newpage

\section{Executive assessment and evidence boundary}
\label{sec:assessment}

\subsection{What should change first}

\textbf{N01: validate the two operands before combining them.}
The highest-priority change is in \code{corePerturbationConstruct}, in
\code{src/Kernel/CorePerturbation.wl}. The constructor tests
\code{FreeQ[core + perturbation, y]}, not
\code{FreeQ[{core, perturbation}, y]}. Its parser then keeps separate core and
perturbation rows, extracts an offset from the core, constructs an exact core
inverse, and computes a marker expansion with a fixed-data majorant.
Those operations depend on facts about the operands that the sum cannot
preserve. The public error message already requires no target symbol in the
forward data; the implementation enforces a weaker condition.\cite{core}

This is not merely an overly optimistic numerical estimate. For the witness in
\cref{sec:counterexample}, the error divided by the stored remainder scale tends
to infinity. No finite hidden big-$O$ constant can repair the result. Its source
and target domains are nonempty and contain the entire required eventual ray.
The bad result is also not dependent on a difficult special-function identity.

\textbf{E01: implement a specific rational certificate family for the Lerch
slice $s=1$.} The repository already has proposals for Lerch remainder bounds and
exact moments. This report does not repeat either request. Its additional
contribution is the explicit two-measure Gauss/Radau construction in
\cref{sec:lerch}: a nested rational interval, a positive denominator formula,
closed error caps, and a linear number of exact recurrence steps. It supplies a
bounded implementation that can be tested independently of either symbolic
kernel.\cite{r49,r50}

\subsection{What actually ran}

\begin{longtable}{@{}p{0.27\textwidth}p{0.65\textwidth}@{}}
\toprule
Evidence class & Actual scope\\\midrule\endhead
Repository reads & Commit-pinned GitHub connector reads of the modules and
review material described below. No local checkout was obtained.\\
Mathematical analysis & Exact all-depth counterexample and fixed-data controls;
proofs of the rational certificate construction and its error identities.\\
Wolfram 15 isolation & A transcribed evaluation of the upstream marker-term
recurrence, the relevant cancellation and realness predicates, and five
coefficients; a separate evaluation of the proposed rational recurrence for
four orders, with four native Lerch comparisons.\\
Python execution & 24 unittest methods, zero failures and zero errors in the
final run. The methods include exact-rational models, independent defining-sum
and moment oracles, and synthetic patch-emitter fixtures.\\
Not executed & The complete original package, the patched package, the supplied
MUnit regressions, the complete supplied WL wrapper file, and all Mathics runs.\\
\bottomrule
\end{longtable}

The available remote kernel reported
\code{15.0.0 for Linux x86 (64-bit) (May 6, 2026)}. Package-loading requests
failed at the service boundary. Container downloads also failed, and no local
Wolfram or Mathics executable was available. These are environment limitations,
not package test failures. Isolated native calculations are consequently marked
as such throughout this article. In particular, the full public N01 call is a
\emph{source-predicted manifestation}, not a fabricated native transcript.

The final Python count must not be added to any historical acceptance count in
the repository. Likewise, the 144 containment subcases in one test method are
not 144 extra unittest methods. The machine-readable scope record and test log
are included in \code{evidence/}.

\subsection{Review coverage}

The audit concentrated on semantic admission, coordinate transformations,
remainder construction, observable operations, exact arithmetic, and portable
proof boundaries. It read the complete exact-core module and selected portions
of the general core, the source and inverse coordinate modules, ordinary and
reciprocal-logarithmic operations, Fourier coefficients, Dirichlet special
functions, native special-function admission, inverse syntax and branches,
certificates, and the Mathics assumption prover. The six wave indexes and the
maintained status register were used to screen candidate findings. Relevant
original sections of reports 49 and 50 were read for the Lerch proposal.

This is a bounded source review, not a statement that every source line, every
retained article, or every test has been re-executed. Unread portions and
unsupported runtimes are not silently converted into coverage. The accompanying
\code{evidence/scope.json} and novelty ledger make that distinction explicit.

\section{Novelty relative to the existing review corpus}
\label{sec:novelty}

The pinned review index retains 45 packages across six waves. Its own accounting
separates report entries from distinct obligations and records overlaps rather
than treating every report-local identifier as a separate bug. That is the
appropriate baseline for an incremental review.\cite{reviewindex,status}

\subsection{N01 is a different admission defect}

The closest existing work is the target-dependent \code{SourceShift} finding
in reports 46 and 50, the fixed-parameter scope principle recorded under C15,
and report 43's source-symbol leakage through an exact core inverse. N01 is
not another example of one of those call paths.\cite{wave5,wave6,status}

\begin{longtable}{@{}p{0.26\textwidth}p{0.66\textwidth}@{}}
\toprule
Comparison & Additional obligation established here\\\midrule\endhead
Reports 46/50 & N01 affects \code{AsymptoticCoreInverse}, not the exponential
constructor's \code{SourceShift} option. The distinguishing mechanism is an
aggregate dependency test that loses information through cancellation.\\
C15 & Parameter capture during composition and admission of two independently
varying operands are different operations. The repair here belongs before the
core model is built, not in substitution or composition.\\
Report 43 & The eliminated source symbol is not retained in N01's inverse.
The target has entered the model through the separately stored core and
perturbation. An output-only source-freedom check would pass this witness.\\
\bottomrule
\end{longtable}

The exponential-core sibling is also an instructive local control: its variable
admission uses a list containing the two operands. This is a precedent for the
minimal repair, not a new request to redesign every chart abstraction.
\cite{expcore}

\subsection{E01 is an algorithm, not a repeated roadmap request}

Report 49 already develops signed mean--variance information for the Lerch
remainder. Report 50 supplies an exact Eulerian moment table and discusses
removing special-function dependence from those moments. The present proposal
uses the same defining positive measure but computes a different object:
two rational resolvent approximants of arbitrary finite order, evaluated by
closed three-term recurrences.\cite{r49,r50}

The added obligations are the two shifted Meixner parameterizations, the
positive-denominator formula, the exact error identities and caps, the nested
interval proof, and a recurrence implementation requiring $O(n)$ arithmetic
steps for a single order-$n$ interval. Order one includes a familiar mean-based
lower bound and is used as a control, not advertised as a new result. The useful
additional family is the higher-order hierarchy and its constructive certificate.

The underlying orthogonal-polynomial and quadrature ideas are classical.
The contribution is a concrete, proved and tested application to this
repository's bounded Lerch certificate lane, not a claim of mathematical priority.

\section{N01: cancellation defeats fixed-operand validation}
\label{sec:counterexample}

\subsection{The source path}

At the constructor boundary the current source contains:
\begin{lstlisting}
validateInput[core + perturbation, limit];
If[x === y || ! FreeQ[core + perturbation, y],
  fail["InvalidVariables",
    "Use distinct source and target symbols, with no target symbol in the forward data."]];
\end{lstlisting}
The public overload separately restricts the coordinate slots to symbols.
The problem is therefore not a missing symbol pattern in this particular
argument form. It is the expression on which dependency is tested.\cite{core}

The subsequent data flow is structurally different:
\begin{lstlisting}
f0 = Simplify[core /. x -> coord["Substitution"], ...];
rr = Simplify[perturbation /. x -> coord["Substitution"], ...];
rows = parseFinite[f0, ...];
remainderRows = parseFinite[rr, ...];
offset = Total[Cases[rows, {0,c_} :> Coefficient[c,ell,0]]];
...
v = y - model["Offset"];
u0 = (v/a)^(1/p);        (* monomial exact core *)
\end{lstlisting}
Here the parser and inverse builder genuinely need the two operands, not merely
their sum. The model records their separate finite supports, the perturbation's
relative source-power gap, and the core offset.\cite{core}

For observable power one at a finite positive-side endpoint, the omitted-tail
pair is constructed as
\begin{lstlisting}
{canon[1 + (depth + 1) delta], (depth + 1) degree}
\end{lstlisting}
when the perturbation is nonzero. The stored majorant explicitly quantifies over
\emph{fixed data}: it asserts existence of finite constants and a threshold for
an expansion in a small source-coordinate scale.\cite{core}

\subsection{A public witness with a trivial true inverse}

Take the target $y$ positive and eventually large, and use
\begin{lstlisting}
AsymptoticCoreInverse[
  1/x + Abs[y]/2,
  -Abs[y]/2,
  {x, 0}, {y, n}
]
\end{lstlisting}
with a fixed nonnegative integer depth $n$. This call is deliberately written
with \code{Abs[y]} rather than a bare $y$ in the coefficient. Native Wolfram
arithmetic can prove \code{Abs[y]} real without a separate assumption about $y$.
The constructor prohibits target-dependent assumptions, so using an unproved-real
bare coefficient would obscure the actual validation defect.

The total forward expression is
\[
 \left(\frac1x+\frac{|y|}{2}\right)-\frac{|y|}{2}=\frac1x.
\]
Consequently the existing \code{FreeQ} test sees no target. On the positive
source chart, both operands are finite power-log expressions. Their source
powers are $-1$ for the leading core term and $0$ for the perturbation. The
computed relative power gap is therefore $\delta=1$, and the logarithmic
degree is zero. The leading amplitude is one. Nothing in those source-power
checks detects that the coefficient at source power zero grows with the target.

On $y>0$, put $A=y/2$. The exact core and perturbation are
\[
 f_0(u)=u^{-1}+A,\qquad r(u)=-A,
 \qquad u_0=\frac1{y-A}=\frac2y.
\]
The stored target-domain ingredients are satisfied for all sufficiently large
$y$: $y-A>0$, $u_0\in\R$, and $0<u_0<1/e$ whenever $y>2e$.
Thus the counterexample does not rely on evaluating outside the selected domain.

\subsection{The marker equation exposes the error}

Introduce the marker $\lambda$ exactly as the algorithm does:
\[
 f_0(u)+\lambda r(u)=y.
\]
For an arbitrary value $A=A(y)$ independent of the source coordinate,
\begin{equation}
 u(\lambda)=\frac1{y-A+\lambda A}
 =\frac{u_0}{1+\lambda A u_0}.
 \label{eq:marker}
\end{equation}
The degree-$j$ marker coefficient is therefore
\begin{equation}
 t_j=(-A)^j u_0^{j+1},\qquad j\geq0.
 \label{eq:marker-coefficient}
\end{equation}
This agrees with the source recurrence based on $r^j/f_0'$ and repeated
application of $(1/f_0')\,\mathrm d/\mathrm du$. It also follows directly by
expanding the rational function in \cref{eq:marker}; no numerical inference is
needed.

For $A=y/2$ one has $A u_0=1$. The finite approximation is
\begin{equation}
 S_n=\sum_{j=0}^{n}t_j
 =\begin{cases}2/y,&n\text{ even},\\0,&n\text{ odd}.\end{cases}
 \label{eq:alternating}
\end{equation}
The true inverse is $u(1)=1/y$, so the absolute error is $1/y$ at every depth.
But the source constructs the remainder pair $(n+2,0)$ and scale
$u_0^{n+2}=(2/y)^{n+2}$. Their ratio is
\begin{equation}
 \frac{|S_n-1/y|}{u_0^{n+2}}
 =\frac{y^{n+1}}{2^{n+2}}\longrightarrow\infty.
 \label{eq:false-scale}
\end{equation}
This refutes the asymptotic remainder at every fixed marker depth.

\begin{theorem}[All-depth cancellation witness]
For each integer $n\geq0$, the source path described above associates a remainder
scale of order $y^{-n-2}$ with an approximation whose error is exactly $y^{-1}$
on an eventual positive target ray. Its claimed big-$O$ statement is false.
\end{theorem}
\begin{proof}
The source-power gap and observable power give $(n+2,0)$. Equations
\eqref{eq:alternating} and \eqref{eq:false-scale} give the finite approximation,
exact error, and unbounded normalized error. The domain check above establishes
that the ray is admitted by the stored domain formula.
\end{proof}

\subsection{The failure survives a genuinely small perturbation ratio}

The previous example has marker ratio $A u_0=1$. A repair that merely checks
whether a marker ratio tends to zero would still miss the more fundamental
fixed-data failure. Set instead
\[
 A(y)=\sqrt y\quad(y>0),
\]
represented in the call by \code{Sqrt[Abs[y]]}. Now
\[
 A u_0=\frac{\sqrt y}{y-\sqrt y}\longrightarrow0.
\]
The marker expansion is geometrically convergent at $\lambda=1$ for all large
enough $y$, and its truncations do improve. Nevertheless they do not improve at
the rate stated in the stored source-power remainder.

\begin{proposition}[General moving-offset identity]
Let $A(y)\geq0$, $y>A(y)$, and $u_0=(y-A(y))^{-1}$. For the degree-$n$ finite
sum from \cref{eq:marker-coefficient},
\begin{align}
 S_n-\frac1y&=-\frac{(-A u_0)^{n+1}}{y},\label{eq:general-error}\\
 \frac{|S_n-1/y|}{u_0^{n+2}}
 &= A(y)^{n+1}\left(1-\frac{A(y)}y\right).
 \label{eq:general-ratio}
\end{align}
If $A(y)=y^\sigma$ with $0<\sigma<1$, the normalized error diverges like
$y^{\sigma(n+1)}$, although $A u_0\to0$.
\end{proposition}
\begin{proof}
Apply the finite geometric-sum identity to \cref{eq:marker}. Since
$u_0/(1+A u_0)=1/y$, the finite remainder has exactly the form
\eqref{eq:general-error}. Dividing its magnitude by $u_0^{n+2}$ gives
\eqref{eq:general-ratio}. The final assertion follows from $A/y\to0$.
\end{proof}

In particular, $A=\sqrt y$ gives an actual error of order
$y^{-(n+3)/2}$, whereas the reported scale is of order $y^{-(n+2)}$.
This second family rules out dismissing the defect as only a borderline
convergence example. The constants in the claimed source-coordinate theorem
are being allowed to depend on the variable that is tending to its limit.

\subsection{The valid fixed-data control}

If $A$ is a genuine fixed real constant, the same algebra yields a bounded
normalized error:
\[
 \frac{|S_n-1/y|}{u_0^{n+2}}
 =|A|^{n+1}\left(1-\frac Ay\right)
\]
on a sufficiently large positive ray. For $y\geq2|A|$ this is at most
$\tfrac32|A|^{n+1}$. A looser factor two is used in the independent tests.
The $A=0$ control is exact and follows the constructor's empty-perturbation path.

The correct fix must therefore distinguish fixed cancellations from moving
cancellations. Rejecting every decomposition in which a constant cancels would
be unnecessary and would remove mathematically valid inputs.

\subsection{Native isolation: what is confirmed}

The successful isolated Wolfram 15 evaluation used the upstream
\code{corePerturbationTerm} body, changing only resource-failure plumbing to a
local throw in a branch that was not reached. Its observed values included:
\begin{lstlisting}
Total forward expression:              1/x
FreeQ[core + perturbation, y]:         True
FullSimplify[Element[perturbation,Reals]]: True
FullSimplify[u0, y>0]:                 2/y
Marker coefficients, degrees 1..5:
  {-2/y, 2/y, -2/y, 2/y, -2/y}
Finite sums, depths 0..5:
  {2/y, 0, 2/y, 0, 2/y, 0}
\end{lstlisting}
This confirms the cancellation gate, the native coefficient-realness control,
and the recurrence mechanics. It does \emph{not} establish that the entire
constructor or the complete patched package was executed. The integration probe
included in this archive is the next required acceptance step, not evidence
that it already ran.

\section{Repair, regression design, and dual-runtime implications}
\label{sec:repair}

\subsection{Minimal semantic change}

The essential repair is to replace the dependency check by
\begin{lstlisting}
If[x === y || ! FreeQ[{core, perturbation}, y],
  fail["InvalidVariables",
    "Use distinct source and target symbols, with no target symbol in either core or perturbation."]];
\end{lstlisting}
The list preserves which expression supplied the data. Ordinary evaluation
inside each operand may still remove a mathematically irrelevant occurrence;
the important point is that addition must not erase a dependency \emph{between}
two operands that the algorithm keeps separate. \code{FreeQ} is a structural
query on its evaluated expression; it cannot recover information already erased
by prior evaluation.\cite{freeq}

The companion change
\begin{lstlisting}
validateInput[{core, perturbation}, limit];
\end{lstlisting}
keeps the exact-input and nonfinite-input checks on the same semantic unit.
For example, inexact constants can cancel in the sum even though the core model
still receives them separately. This companion change is an explicit strengthening
of operand admission, not a separate counted bug and not a claim that every such
input previously produced a wrong answer.

The patch retains the existing failure tag \code{InvalidVariables}. The total
forward function may still be stored as \code{core + perturbation}; it is the
right expression for the defining equation after both operands have passed
admission. There is no need to change the valid fixed-data marker formula,
change the result head, or reclassify this object as an exponent-sorted jet.

\subsection{Why an output-only repair is inadequate}

Three apparently reasonable postconditions would all miss the first witness:
the total forward expression is target-free, the exact core inverse is free of
the eliminated source symbol, and every displayed coefficient is real on the
target ray. The obstruction belongs to the model's quantifiers, not to the
syntactic cleanliness of the final expression.

Likewise, multiplying the first omitted term by a guessed numerical constant
would not repair the declared remainder theorem. The all-depth ratio proof
shows that a constant independent of the target does not exist for the stored
scale. A general moving-parameter theory would need its own asymptotic grading,
but the existing public contract calls for rejecting this input before invoking
a fixed-data theorem. That much narrower change is the recommended repair.

\subsection{Regression expectations}

The supplied \code{tests/CoreOperandScope.wlt} is a desired-contract MUnit file.
It is unrun. It requires both moving-coefficient families to return the
structured variable-admission failure, keeps a fixed cancellation as a
\code{GeneralizedSeries}, checks the exact zero-perturbation control, and checks
that inexact operands do not disappear before the exact-input boundary.

The separate \code{ProbeCoreOperandScope.wl} is suitable for observational use
in either session after loading the package. It prints the runtime and explicitly
distinguishes a series result from a refusal. For a series it reads the actual
\code{CoreInverse} fields \code{Expression}, \code{RemainderPower},
\code{RemainderScaleExpression}, and \code{TargetDomain}. It does not reuse an
ordinary-forward-field assumption or count an unevaluated numerical guard as a
passing comparison.

A useful native acceptance sequence is: run the probe on the original pin;
apply the two-line semantic repair to a separate working copy; rebuild the
standalone using the repository's maintained procedure; rerun the probe and
MUnit file; then execute existing exact-core controls. The patch emitter does
not perform these steps or claim their results.

\subsection{Mathics: common source is not common observed behavior}

The offending dependency check is shared source, so it should be corrected
independently of which runtime first reproduces the bad public call. However,
the particular \code{Abs[y]} witness relies on a realness fact that native
Wolfram can discharge. The Mathics assumption adapter has a more conservative
recursive realness prover; it may refuse the expression before reaching the
same public result. Reading common source is not evidence that this public
witness succeeds identically in Mathics.\cite{mathicsassumptions}

The precise claims are therefore: the shared guard loses operand information;
the native isolated realness and recurrence controls pass; the full Wolfram
manifestation follows the inspected source path but was not integration-run;
and Mathics public reachability remains unmeasured. The proposed list-based
admission itself uses elementary operations already present in the sibling
implementation and does not introduce a new special-function dependency.

The Lerch prototype in \cref{sec:lerch} is intentionally separate. It uses exact
rational operations, finite loops, lists, and a small association wrapper.
That is a design for portability, not proof of runtime compatibility. A fresh
Mathics run of the supplied file is still required.

\section{E01: a rational Lerch certificate hierarchy}
\label{sec:lerch}

\subsection{The bounded target}

For fixed $0<q<1$ and $a>0$, define
\begin{equation}
 F_q(a)=\PhiL(q,1,a)=\sum_{k=0}^{\infty}\frac{q^k}{a+k}.
 \label{eq:lerch}
\end{equation}
This is the $s=1$ slice of the defining Lerch series.\cite{dlmflerch}
The proposed primitive returns an exact interval for the \emph{function value},
not a replacement remainder for an arbitrary generalized-series object.
Its first implementation accepts rational $q,a$ and a positive integer order.
It treats $q=0$ separately, where $F_0(a)=1/a$ exactly.

The restriction is useful. A positive discrete measure is available, all moments
are finite, and both denominator families have explicit positive evaluations
on the negative real axis. No nonprincipal branches or numerical sign decisions
enter the certificate. Extending to $q<0$, complex parameters, or arbitrary $s$
is outside this prototype.

\subsection{Two positive measures, not a moment matrix}

Let
\[
 \mu=\sum_{k\geq0}q^k\delta_k,
 \qquad \mu_1=\sum_{k\geq0}kq^k\delta_k.
\]
Their masses are
\begin{equation}
 m_0=\frac1{1-q},\qquad m_1=\frac{q}{(1-q)^2}.
 \label{eq:masses}
\end{equation}
Then $F_q(a)=\int(a+t)^{-1}\dd\mu(t)$. The weighted transform satisfies
\begin{equation}
 H_q(a)=\int\frac1{a+t}\dd\mu_1(t)
       =m_0-aF_q(a).
 \label{eq:weighted}
\end{equation}
The first measure has support in $[0,\infty)$; the second has support in
$[1,\infty)$ after its zero weight at $k=0$ is discarded. This one-unit shift
will improve its elementary error cap.

The first measure is the Meixner weight with parameter $\beta=1$. Writing
$k=t+1$ in the second gives $q(t+1)q^t$, namely $q$ times the Meixner weight
with $\beta=2$, translated by one. The defining Meixner orthogonality and
finite hypergeometric representation are recorded in DLMF
\S\S18.19--18.20.\cite{dlmfmeixner,dlmfexplicit}

\subsection{A general resolvent identity}

Let $P_n$ be the monic degree-$n$ polynomial orthogonal for a positive measure
$\nu$ supported in $[b,\infty)$, where $b\geq0$. Define
\begin{equation}
 G_n^{\nu}(a)=\frac1{P_n(-a)}
 \int\frac{P_n(-a)-P_n(t)}{a+t}\dd\nu(t).
 \label{eq:gauss}
\end{equation}
The quotient in the integral is a polynomial of degree $n-1$, so this is a
rational function computable from finite polynomial information. The formula
also identifies the usual Gaussian resolvent approximant without requiring
its quadrature nodes.

\begin{theorem}[Positive exact error identity]
Assume $a>0$ and $P_n(-a)\ne0$. Then
\begin{equation}
 \int\frac1{a+t}\dd\nu(t)-G_n^{\nu}(a)
 =\frac1{P_n(-a)^2}\int\frac{P_n(t)^2}{a+t}\dd\nu(t).
 \label{eq:gauss-error}
\end{equation}
In particular the error is positive for a nondegenerate measure, and with
$h_n^{\nu}=\int P_n^2\dd\nu$ it is at most
\begin{equation}
 \frac{h_n^{\nu}}{(a+b)P_n(-a)^2}.
 \label{eq:norm-bound}
\end{equation}
\end{theorem}
\begin{proof}
The polynomial $(P_n(t)-P_n(-a))/(t+a)$ has degree $n-1$. Orthogonality gives
\[
 0=\int P_n(t)\frac{P_n(t)-P_n(-a)}{a+t}\dd\nu(t).
\]
Rearranging yields
$\int P_n(t)/(a+t)\dd\nu(t)=
 P_n(-a)^{-1}\int P_n(t)^2/(a+t)\dd\nu(t)$.
Subtract \cref{eq:gauss} from the original transform and substitute this
identity. Positivity follows from the integrand. Finally $a+t\geq a+b$
gives \cref{eq:norm-bound}.
\end{proof}

\subsection{The lower and upper endpoints}

Use $G_n=G_n^{\mu}$ and $J_n=G_n^{\mu_1}$, and set
\begin{equation}
 L_n(a)=G_n(a),\qquad
 U_n(a)=\frac{m_0-J_n(a)}a.
 \label{eq:bounds}
\end{equation}
By \cref{eq:weighted,eq:gauss-error},
\begin{equation}
 L_n(a)<F_q(a)<U_n(a)\qquad(0<q<1,\ a>0).
 \label{eq:strict-enclosure}
\end{equation}
The upper endpoint is the Gauss/Radau-type bound obtained by applying the
Gaussian lower-bound construction to the size-biased measure. Both endpoints
are rational when $q,a$ are rational.

This avoids a common integration mistake: an accurate lower approximation to
$H_q$ becomes an \emph{upper} approximation to $F_q$ because of the minus sign
in \cref{eq:weighted}. The implementation keeps the two measures separate and
its independent oracle checks both orientations.

\subsection{Closed recurrence coefficients}

For a Meixner parameter $\beta>0$ translated by $\delta\geq0$, the monic
three-term recurrence has coefficients
\begin{align}
 b_j&=\delta+\frac{(1+q)j+\beta q}{1-q},\label{eq:diagonal}\\
 \lambda_j&=\frac{qj(j+\beta-1)}{(1-q)^2},\qquad j\geq1.
 \label{eq:subdiagonal}
\end{align}
They follow by substituting the finite Meixner representation into the monic
recurrence, or by matching its polynomial coefficients. The independent tests
also reconstruct the finite polynomials and their exact moments, rather than
using this recurrence as their sole oracle.

For total mass $m$, initialize
\[
 D_0=1,\quad D_1=a+b_0,\qquad N_0=0,\quad N_1=m.
\]
For $j=1,\ldots,n-1$, compute
\begin{align}
 D_{j+1}&=(a+b_j)D_j-\lambda_jD_{j-1},\\
 N_{j+1}&=(a+b_j)N_j-\lambda_jN_{j-1}.
 \label{eq:recurrence}
\end{align}
The rational approximant is $N_n/D_n$. The denominator is
$D_n=(-1)^nP_n(-a)$. For $L_n$, use $(\beta,\delta,m)=(1,0,m_0)$;
for $J_n$, use $(2,1,m_1)$.

Only two consecutive numerator and denominator values are needed. For a
single requested order $n$, two recurrences cost $O(n)$ exact arithmetic
operations and $O(1)$ live scalar slots. This is not an $O(n)$ bit-complexity
claim: the numerators and denominators grow, and rational reductions have
nonconstant cost. The proposal nevertheless removes moment-matrix construction,
linear-system solution, root isolation and special-function evaluation from
this certificate's production path.

\subsection{Denominator positivity without approximate roots}

A subtraction recurrence deserves a denominator proof even when sample outputs
look positive. The finite Meixner formula gives
\begin{equation}
 D_n(a;\beta,\delta)=
 (\beta)_n\left(\frac q{1-q}\right)^n
 \sum_{j=0}^{n}\binom nj
 \frac{(a+\delta)_j}{(\beta)_j}
 \left(\frac{1-q}{q}\right)^j.
 \label{eq:positive-denominator}
\end{equation}
Every summand is positive in the stated parameter domain. Thus $D_n>0$
exactly, without computing a root or rounding a sign. The production code
additionally checks its positive-denominator invariant to catch an implementation
error; the mathematical guarantee is \cref{eq:positive-denominator}, not that
runtime check.

For the two required parameter choices, all factors are rational at rational
$a,q$. Although \cref{eq:positive-denominator} is also an evaluation algorithm,
the recurrence is used in production; the finite sum is an independent exact
check and a convenient proof certificate.

\subsection{Nested bounds}

Let $W_n=N_nD_{n-1}-N_{n-1}D_n$. Equation \eqref{eq:recurrence} gives
\[
 W_{n+1}=\lambda_nW_n,\qquad W_1=m.
\]
Consequently
\begin{equation}
 \frac{N_{n+1}}{D_{n+1}}-\frac{N_n}{D_n}
 =\frac{m\lambda_1\cdots\lambda_n}{D_nD_{n+1}}>0.
 \label{eq:nesting}
\end{equation}
The lower endpoints increase. The Gaussian approximants to $H_q$ also increase,
so \cref{eq:weighted} makes the upper endpoints decrease:
\[
 L_n<L_{n+1}<F_q<U_{n+1}<U_n.
\]
Thus increasing the order retains the previous enclosure while shrinking its
width. This is a property proved for the exact construction, not an empirical
rule inferred from the numerical examples.

\subsection{Explicit error caps and asymptotic order}

The monic squared norms are
\begin{align}
 h_n&=\frac{(n!)^2q^n}{(1-q)^{2n+1}},\label{eq:norm0}\\
 h_n^{(1)}&=\frac{n!(n+1)!q^{n+1}}{(1-q)^{2n+2}}.
 \label{eq:norm1}
\end{align}
These follow from the Meixner normalization, or from $h_n=m\prod_{j=1}^n
\lambda_j$ with the two parameter choices.\cite{dlmfmeixner}
Writing $D_n^{(1)}$ for the denominator of $J_n$ gives the separately valid caps
\begin{align}
 0<F_q-L_n&\leq \frac{h_n}{aD_n^2},\label{eq:lower-cap}\\
 0<U_n-F_q&\leq
 \frac{h_n^{(1)}}{a(a+1)(D_n^{(1)})^2}.
 \label{eq:upper-cap}
\end{align}
The factor $a+1$ uses the actual support of $\mu_1$; replacing it by $a$
would remain sound but lose information.

At fixed $q,n$, the exact error identities also give
\begin{align}
 F_q-L_n&=h_n a^{-2n-1}+O(a^{-2n-2}),\\
 U_n-F_q&=h_n^{(1)}a^{-2n-2}+O(a^{-2n-3}).
\end{align}
To justify these statements, $D_n$ is monic in $a$, and the difference between
$\int P_n(t)^2/(a+t)\dd\nu$ and $h_n/a$ is bounded in magnitude by
$a^{-2}\int tP_n(t)^2\dd\nu$. The required moment is finite. No uniformity as
$q\to1$ is claimed; the displayed norms explicitly reveal the deterioration
near that endpoint.

\subsection{Concrete exact example}

For $q=1/2$, $a=10$, the first four intervals are:
\begin{center}
\begin{tabular}{@{}rrrr@{}}\toprule
$n$ & $L_n$ & $U_n$ & $U_n-L_n$\\\midrule
1 & $2/11$ & $12/65$ & $2/715$\\
2 & $7/38$ & $47/255$ & $1/9690$\\
3 & $115/624$ & $857/4650$ & $1/161200$\\
4 & $1087/5898$ & $8863/48090$ & $4/7878745$\\
\bottomrule\end{tabular}
\end{center}
The corresponding width displays are approximately
$2.80\cdot10^{-3}$, $1.03\cdot10^{-4}$, $6.20\cdot10^{-6}$, and
$5.08\cdot10^{-7}$. The fractions, not these decimal displays, are the
certificates. At order two the additional caps are
\[
 F_{1/2}(10)-\frac7{38}\leq\frac1{7220},\qquad
 \frac{47}{255}-F_{1/2}(10)\leq\frac1{47685}.
\]
The isolated Wolfram 15 recurrence evaluation returned precisely the four
rational endpoint pairs and their error caps. Four native comparisons with
\code{LerchPhi[1/2,1,10]} returned \code{True}. The independent exact oracle
provides a separate check not based on that built-in function.

\section{Implementation contracts and tests for E01}
\label{sec:implementation}

\subsection{A small API with explicit semantics}

The Python entry point is
\begin{lstlisting}
lerch_bounds(q: int | Fraction,
             a: int | Fraction,
             order: int) -> LerchBounds
\end{lstlisting}
It returns exact endpoints, their width, two separate error caps, the requested
order, the input values, and a method identifier. Approximate inputs, strings,
booleans used as numbers, invalid domains, and noninteger orders are rejected.
The prototype caps the order at 64 as an explicit resource policy; the theorem
is not limited to 64. Its degenerate $q=0$ branch returns an exact point interval.

The WL counterpart is \code{ReviewLerch`LerchRationalBounds[q,a,n]} in a separate
context. It returns a small association or a structured failure. It does not
install definitions on \code{System`} symbols, alter the repository's private
helpers, or modify \code{LerchPhi}. The supplied wrapper file has not been run
as a complete file in either runtime; the native observation executes the same
arithmetic recurrence body with an isolation wrapper.

\subsection{Why value bounds must not become series metadata accidentally}

An enclosure of $F_q(a)$ is not automatically an error bound for an arbitrary
stored truncated expression $T(a)$. At a given exact parameter value, the valid
transport is
\[
 F_q(a)-T(a)\in[L_n(a)-T(a),\ U_n(a)-T(a)].
\]
An absolute error bound is then the maximum of the absolute values of these two
endpoints, provided $T(a)$ itself has been evaluated exactly or enclosed
independently. A decimal approximation to $T(a)$ is not sufficient for that
step. The new primitive must not silently declare a derivative contract, a
uniform parameter-domain theorem, or an inverse-branch certificate.

A first integration should therefore remain an opt-in value-certificate helper.
It can attach its exact input point, theorem family, order, endpoints and
arithmetic domain to a result. Promotion into more general remainder metadata
should be a separate, explicitly tested transport operation. This is a concrete
contract for this new family, not a replacement proposal for the repository's
already-recorded result-schema work.

\subsection{Independent validation paths}

The production recurrence is checked in three different ways.
First, the defining sum gives the entirely rational enclosure
\begin{equation}
 S_K\leq F_q(a)\leq S_K+
 \frac{q^K}{(a+K)(1-q)},\qquad
 S_K=\sum_{k=0}^{K-1}\frac{q^k}{a+k}.
 \label{eq:sum-oracle}
\end{equation}
The test increases $K$ until this oracle interval lies inside the proposed
interval and is fine enough to verify both explicit error caps.

Second, test-only code constructs Meixner polynomials from their terminating
finite formula. It forms the polynomial quotient in \cref{eq:gauss} and
integrates it using exact factorial moments. Those moments are obtained from
\[
 \sum_{k\geq0}\frac{(\beta)_k}{k!}q^k
 k^{\underline j}
 =\frac{(\beta)_j q^j}{(1-q)^{\beta+j}}.
\]
A Stirling-number conversion supplies ordinary moments. This is a different
construction from the production three-term recurrence.

Third, the tests verify orthogonality and squared norms exactly, compare the
positive finite denominator sum with the recurrence, and check interval nesting.
These are algebraic identities, not floating-point residual thresholds.

\begin{center}
\begin{tabular}{@{}lr@{}}\toprule
Selected subcase population & Count\\\midrule
Defining-sum containment and error-cap cases & 144\\
Exact orthogonality identities & 224\\
Moment/quotient endpoint comparisons & 144\\
Adjacent-order nesting comparisons & 264\\
Moving-parameter inverse cases & 66\\
Largest defining-sum cutoff needed & 512\\
\bottomrule\end{tabular}
\end{center}
These populations belong to the 24 unittest methods; they are not additive
acceptance counts. The parameter grid for containment is
$q\in\{1/5,1/2,3/4,9/10\}$,
$a\in\{1/2,1,2,10,100,1000\}$ and $n=1,\ldots,6$.

\subsection{A test-harness correction, preserved in the evidence}

The first Python pass had two subcase failures. The oracle-refinement loop
initially stopped as soon as its interval lay inside $[L_n,U_n]$, then tried to
verify the tighter individual error caps with the same insufficiently refined
oracle. The endpoint-enclosure tests had succeeded; the oracle did not yet have
enough resolution to establish the separate cap inequalities.

The correction was to require all four inequalities in the loop's stopping
condition. Neither the recurrence nor the mathematical caps were changed.
The final pass has 24 successful methods and no failures. Both first-pass and
final records are included so that the improvement in the test harness is not
misrepresented as an unexplained change in mathematical output.

\subsection{Performance claims that are and are not supported}

For each measure the implementation performs $n-1$ recurrence steps after
initialization, and it updates a norm product at the same time. It keeps only
two consecutive values of each recurrence. Those structural work bounds follow
from the source and can be inspected directly.

No Wolfram-versus-Mathics timing comparison was run. There is no claim that this
prototype outperforms the existing package on a measured workload, that order
64 is an optimal cap, or that exact arithmetic prevents all expression growth.
Near $q=1$ or with large rational input heights, numerator growth can dominate.
A production implementation should expose the cost of this particular finite
certificate separately from the approximation order of an unrelated series.

\section{Other inspected candidates not counted as findings}
\label{sec:negative}

A careful incremental audit needs to discard attractive but unsupported claims.
One examined path was the terminal branch of the private exponent comparator,
which can infer a positive ordering from a simplified negation of a negative
comparison. In isolation, $\neg(d<0)$ only implies $d\geq0$, not $d>0$.
However, the actual comparator first canonicalizes its exact inputs and checks
for zero. A public counterexample must survive that earlier stage and satisfy
the caller's admitted input domain; a purely symbolic private-helper example
would not establish this.\cite{maincore}

Two proposed noncanonical-zero witnesses were rejected during this audit.
A Bessel Wronskian combination and a rational-argument dilogarithm five-term
identity both reduced to zero under the available native \code{FullSimplify}
control. The latter evaluation returned a canonical absolute difference of zero.
That canonical zero is decisive: it prevents the alleged
path through the comparator's final branch. No second public ordering defect or
repair is claimed from those examples.

This matters for the evidence standard. A suspicious implication in a helper,
a failed direct numerical comparison, and a mathematical identity are not enough
when the real control flow has an earlier successful normalization. The article
therefore retains N01 as its one new correctness finding rather than inflating
the count with an unestablished public manifestation.

\section{Integration priorities and reproducibility}
\label{sec:reproduce}

\subsection{Recommended sequence}

The first change should be the operand-preserving N01 admission rule, with the
moving-small family as well as the alternating family in acceptance. The valid
fixed-cancellation and exact-zero controls are essential: they prevent an
unnecessarily broad rejection policy.

The second stage is a focused execution of those requests in each runtime.
A Mathics refusal caused by a separate realness limitation should be recorded
as such, not described as proof that the underlying shared admission is sound.
The generated standalone and modular loading routes should agree after rebuilding.
This report's source patch has not established that agreement.

The rational Lerch primitive should be integrated separately. First run the
supplied WL wrapper on the same exact fixtures as the Python implementation.
Then attach it only to the $s=1$, $0\leq q<1$, $a>0$ value-certificate path.
Broader symbolic parameters or automatic dispatch should wait for a separately
stated admission and transport contract. This sequencing prevents a small
correctness repair from being entangled with an unrelated algorithmic extension.

\subsection{Archive layout}
\begin{lstlisting}
article/article.tex
article/article.pdf
code/emit_core_patch.py
code/lerch_rational.py
code/LerchRationalBounds.wl
tests/test_review.py
tests/run_tests.py
tests/CoreOperandScope.wlt
tests/ProbeCoreOperandScope.wl
evidence/...
README.md
build.sh
LICENSE
\end{lstlisting}
The archive contains no checksum files and no complete copy of the upstream
repository. The patch emitter reads a local checkout supplied by the reader and
emits a unified diff without modifying source. It requires each exact source
anchor to occur once; missing, duplicated, or already-patched anchors are
rejected. Its tests use synthetic fixtures and do not prove that a full checkout
was patched successfully.

\subsection{Commands}

From the archive root, the independent tests require only a standard Python
installation supporting the supplied syntax:
\begin{lstlisting}
python tests/run_tests.py
python code/lerch_rational.py 1/2 10 4
python code/emit_core_patch.py /path/to/Asymptotic --output core-operands.patch
\end{lstlisting}
The test runner regenerates the final evidence files; preserve the supplied
records elsewhere before overwriting them for a new environment. The initial
historical pass is retained separately.

For a symbolic session, load the files explicitly:
\begin{lstlisting}
Get["/path/to/archive/code/LerchRationalBounds.wl"];
ReviewLerch`LerchRationalBounds[1/2,10,4]

Get["/path/to/Asymptotic/AsymptoticAnalysis.wl"];
Get["/path/to/archive/tests/ProbeCoreOperandScope.wl"];
\end{lstlisting}
Use the repository's own rebuilding instructions after applying a patch to its
modular sources. A generated standalone must not be assumed current merely
because the modular file was edited. This archive does not automate or assert
that repository-specific rebuild.

The article builds with three \code{pdflatex} passes; \code{build.sh} runs them in
a temporary directory and copies only the final PDF into \code{article/}.
No bibliography program, network request, or proprietary font is required.

\section{Conclusion}

The new correctness result is a failure of \emph{operand-level admission}, not
of the marker coefficient formula. Checking the simplified sum establishes a
property of the total forward function but not of the separately retained model
data. The exact counterexample shows why that distinction matters: the defining
function is elementary, its inverse is unambiguous, and yet the stored asymptotic
scale is false at every depth. A genuinely small moving perturbation produces
the same contract failure, so convergence screening alone cannot repair it.

The rational Lerch proposal illustrates a complementary development direction:
identify a narrow mathematical family, derive the certificate including its
denominator and remainder obligations, and implement the arithmetic without
requiring the symbolic kernel to rediscover those facts. Here two explicit
Meixner recurrences provide a nested exact interval and separately bounded
endpoint errors. The independent tests and isolated native controls support that
prototype at their stated scopes. Full package integration and Mathics execution
remain concrete acceptance work, not results claimed by this review.

\appendix
\section{Source and evidence map}

\begin{longtable}{@{}p{0.31\textwidth}p{0.61\textwidth}@{}}
\toprule
Claim & Primary evidence\\\midrule\endhead
N01 sum-only admission & \code{corePerturbationConstruct} in
\code{CorePerturbation.wl}; pinned source in reference \cite{core}.\\
Separate model operands & \code{corePerturbationModel}, including separate
\code{rows}, \code{remainderRows}, and \code{Offset}.\\
Exact inverse and marker terms & \code{corePerturbationAutomaticInverse},
\code{corePerturbationTerm}, and constructor assembly.\\
Remainder theorem quantified over fixed data & Constructor
\code{truncationPair}, \code{scale}, and \code{MajorantContract}.\\
Native recurrence observation & \code{evidence/native-observations.json},
record N01-isolated-recurrence; transcribed selected tool output.\\
Native rational prototype observation & Same file, record E01-isolated-recurrence;
not a complete wrapper-file or package run.\\
Independent models and fixtures & \code{tests/test_review.py},\newline
\code{evidence/python-results.json}, and \code{python-tests.txt}.\\
First test-oracle pass & \code{python-first-pass.json} and
\code{python-first-pass.txt}; two subcase failures before oracle refinement.\\
Exact rational sample intervals & \code{evidence/lerch-exact-samples.json}.
Decimal values in the CSV are display-only.\\
Novelty and omissions & \code{evidence/novelty-ledger.json} and
\code{evidence/scope.json}.\\
\bottomrule
\end{longtable}

\begin{thebibliography}{99}\small
\bibitem{core} V. Reshetnikov, \emph{AsymptoticAnalysis},
\code{src/Kernel/CorePerturbation.wl}, reviewed commit
\code{efa1aeec4845a9c35e140963a0333d0c9ec33b05}.
\url{https://github.com/VladimirReshetnikov/Asymptotic/blob/efa1aeec4845a9c35e140963a0333d0c9ec33b05/src/Kernel/CorePerturbation.wl}.
\bibitem{expcore} Same repository and commit,
\code{src/Kernel/ExponentialCorePerturbation.wl}.
\url{https://github.com/VladimirReshetnikov/Asymptotic/blob/efa1aeec4845a9c35e140963a0333d0c9ec33b05/src/Kernel/ExponentialCorePerturbation.wl}.
\bibitem{maincore} Same repository and commit,
\code{src/Kernel/AsymptoticAnalysis.wl}, exact-real canonicalization and comparison.
\url{https://github.com/VladimirReshetnikov/Asymptotic/blob/efa1aeec4845a9c35e140963a0333d0c9ec33b05/src/Kernel/AsymptoticAnalysis.wl}.
\bibitem{mathicsassumptions} Same repository and commit,
\code{src/Kernel/MathicsAssumptions.wl}, finite-real and sign proof rules.
\url{https://github.com/VladimirReshetnikov/Asymptotic/blob/efa1aeec4845a9c35e140963a0333d0c9ec33b05/src/Kernel/MathicsAssumptions.wl}.
\bibitem{reviewindex} Same repository and commit, review index and six linked wave indexes.
\url{https://github.com/VladimirReshetnikov/Asymptotic/blob/efa1aeec4845a9c35e140963a0333d0c9ec33b05/external-reports/code-review/README.md}.
\bibitem{status} Same repository and commit,
\code{docs/development/CODE_REVIEW_STATUS.md}.
\url{https://github.com/VladimirReshetnikov/Asymptotic/blob/efa1aeec4845a9c35e140963a0333d0c9ec33b05/docs/development/CODE_REVIEW_STATUS.md}.
\bibitem{wave5} Same repository and commit, wave-5 review index.
\url{https://github.com/VladimirReshetnikov/Asymptotic/blob/efa1aeec4845a9c35e140963a0333d0c9ec33b05/external-reports/code-review/wave-5/README.md}.
\bibitem{wave6} Same repository and commit, wave-6 review index.
\url{https://github.com/VladimirReshetnikov/Asymptotic/blob/efa1aeec4845a9c35e140963a0333d0c9ec33b05/external-reports/code-review/wave-6/README.md}.
\bibitem{r49} Retained review 49, \emph{Incremental Contract Audit of
AsymptoticAnalysis}, section on signed Lerch remainder bounds, at the reviewed commit.
\url{https://github.com/VladimirReshetnikov/Asymptotic/blob/efa1aeec4845a9c35e140963a0333d0c9ec33b05/external-reports/code-review/wave-6/code-review-49/article.tex}.
\bibitem{r50} Retained review 50, \emph{Fixed charts, exact cores, rational
moments}, section on Eulerian Lerch moments, at the reviewed commit.
\url{https://github.com/VladimirReshetnikov/Asymptotic/blob/efa1aeec4845a9c35e140963a0333d0c9ec33b05/external-reports/code-review/wave-6/code-review-50/article/article.tex}.
\bibitem{freeq} Wolfram Research, \emph{FreeQ}, Wolfram Language documentation,
accessed 10 September 2026.
\url{https://reference.wolfram.com/language/ref/FreeQ.html}.
\bibitem{dlmflerch} NIST Digital Library of Mathematical Functions,
\S25.14, \emph{Lerch's Transcendent}, accessed 10 September 2026.
\url{https://dlmf.nist.gov/25.14}.
\bibitem{dlmfmeixner} NIST DLMF, \S18.19, \emph{Hahn Class: Definitions},
Meixner weights, norms and leading coefficients, accessed 10 September 2026.
\url{https://dlmf.nist.gov/18.19}.
\bibitem{dlmfexplicit} NIST DLMF, \S18.20, \emph{Hahn Class: Explicit
Representations}, Meixner finite hypergeometric representation, accessed
10 September 2026. \url{https://dlmf.nist.gov/18.20}.
\end{thebibliography}
\end{document}
